\chapter{总结展望}

本报告以统一结构比较了六个地学智能赛道。结果显示，模型选择必须服从任务的数据
形态和样本规模。树模型在储层物性与岩相任务中仍具有稳定优势；SAM2 交叉注意力和
Chronos-2 分别在地震相与 T3 开发实验中取得改进。三维重建则通过改变基础模型的
接入位置取得进展：冻结 OpenMind-MAE 表征不再直接回归孔隙度，而是参与邻域度量，
固定三核集成在开发空间折和同场区历史版本迁移中均低于 PyKrige 的 RMSE。

新增空间诊断进一步说明，点值指标并不足以支撑地学模型。断层结果需要回到地震剖面
核对几何位置，多个物性目标需要检查联合关系，生产预测需要检查时序残差，三维重建
还必须评价方向空间相关。P21 改善了点值误差，但现有条件重建的方向变差仍小于真值；
这两项观察并不冲突，它们分别描述数值预测和细尺度空间结构。

\begin{table}[htbp]
  \centering
  \caption{六个赛道的主要结果与后续实验}
  \label{tab:conclusion}
  \small
  \begin{tabularx}{\textwidth}{>{\raggedright\arraybackslash}p{2.50cm}YY}
    \toprule
    赛道 & 主要结果 & 后续实验 \\
    \midrule
    断层预测 & 二维基线高召回、低精确率 & 建立连续三维开发折 \\
    地震相分割 & 开发集 mIoU 提高 0.1160 & 补同架构随机初始化对照 \\
    储层物性预测 & 树模型留出 $R^2$ 为 0.8891--0.9807 & 校准长尾区间并补新井族 \\
    岩相分类 & MOMENT 略优于随机初始化但弱于 XGBoost & 扩大母井与少数类样本 \\
    甜点评价 & T3 开发 MAE 为 186.57 & 完成同口径留出评价 \\
    三维重建 & P21 开发 RMSE 为 0.027734，历史版本迁移改善 1.4529\% & 冻结 P21，继续检查细尺度空间结构 \\
    \bottomrule
  \end{tabularx}
\end{table}

后续有三项优先工作。第一，为断层预测建立合法的三维开发折，使 SAM-Med3D
候选模型具备公平评价条件。第二，为地震相基础模型分支补充同架构随机权重对照，
并为 Chronos-2 完成同口径时序留出评价；岩相分类已完成预训练归因，应转向扩大
数据与改善少数类学习。第三，冻结三维重建的 P21 默认配置，不再扩展未晋级的
LoRA、Adapter 或残差旁路。当前报告不追加跨数据集测试，后续工作转向检查细尺度
空间结构和不确定性。若未来获得独立新井，再将其作为额外验证，而不回写本次模型选择。

报告结构已经为后续实验预留扩展空间。新增模型时，应在对应赛道的“方法原理”和
“训练策略”中说明变化。新增数据或对照时，应更新“实验设计”。新增结果则进入
“实验结果”，并保持原有指标口径。这样可以持续扩充内容，而不破坏六赛道之间的
可比性。
